\section{Ôn tập chương 8 - Xác suất}
\Opensolutionfile{ans}[ans/ans-OTC8-TN]
\caulc
\begin{ex}%Câu 1.%[1D9N1-1]
	Cặp biến cố A và B được gọi là độc lập nếu
	\choice
	{Việc biến cố này xảy ra hay không ảnh hưởng tới xác suất xảy ra của biển cố kia}
	{\True Việc xảy ra hay không xảy ra của biến cố này không ảnh hưởng tới xác suất xảy ra của biến cố kia}
	{A hoặc B xảy ra}
	{A và B cùng xảy ra}
	\loigiai{
		Dựa vào định nghĩa.}
\end{ex}
\begin{ex}%Câu 2.%[1D9N2-1]
	Cho hai biến cố A và B. Biến cố hợp của A và B là biến cố
	\choice
	{\lq\lq A và B xảy ra\rq\rq }
	{\True \lq\lq A hoặc B xảy ra\rq\rq }
	{\lq\lq A xảy ra\rq\rq }
	{\lq\lq B xảy ra hoặc cả A và B xảy ra\rq\rq }
	\loigiai{
		Dựa vào định nghĩa.}
\end{ex}
\begin{ex}%Câu 3.%[1D9N1-1]
	Cho hai biến cố A và B. Biến cố giao của A và B là biến cố: 
	\choice
	{\True \lq\lq A và B xảy ra\rq\rq }
	{\lq\lq A hoặc B xảy ra\rq\rq }
	{\lq\lq A xảy ra\rq\rq }
	{\lq\lq B xảy ra hoặc cả A và B xảy ra\rq\rq }
	\loigiai{
		Dựa vào định nghĩa.}
\end{ex}
\begin{ex}%Câu 4.%[1D9N2-1]
	Hai bạn Triết và Khang, mỗi người gieo đồng thời một con xúc xắc cân đối, đồng chất. Xét hai biến cố sau:\\
	A: \lq\lq Số chấm xuất hiện trên con xúc xắc bạn Triết gieo là số lẻ.\rq\rq .\\
	B: Số chấm xuất hiện trên con xúc xắc bạn Khang gieo là số chia hết cho 5.\rq\rq .\\
	Chọn câu đúng: 
	\choice
	{\True Hai biến cố A, B là hai biến cố độc lập}
	{Hai biến cố A, B là hai biến cố xung khắc}
	{Biến cố A giao B là $\{1;3;5\}$}
	{Biến cố A hợp B là $\{1;2;3;5\}$}
	\loigiai{
		Dựa vào định nghĩa.}
\end{ex}
\begin{ex}%Câu 5.%[1D9N2-1]
	Một hộp có $12$ chiếc thẻ cùng loại, mỗi thẻ được ghi một trong các số tự nhiên, hai thẻ khác nhau thì ghi hai số khác nhau. Rút ngẫu nhiên 1 chiếc thẻ trong hộp. Xét biến cố A: \rq\rq Số xuất hiện trên thẻ được rút ra là số chia hết cho 3\lq\lq\, và biến cố B: \lq\lq Số xuất hiện trên thẻ được rút ra là số chia hết cho 4\rq\rq.\\
	Biến cố A hợp B được phát biểu như sau: 
	\choice
	{\lq\lq Số xuất hiện trên thẻ là số vừa chia hết cho 3 vừa chia hết cho 4\rq\rq }
	{\True \lq\lq Số xuất hiện trên thẻ là số chia hết cho 3 hoặc chia hết cho 4\rq\rq }
	{\lq\lq Số xuất hiện trên thẻ là số chia hết cho 12\rq\rq }
	{\lq\lq Số xuất hiện trên thẻ là số chia hết cho 7\rq\rq }
	\loigiai{
		Dựa vào định nghĩa.}
\end{ex}
\begin{ex}%Câu 6.%[1D9N1-1]
	Chọn ngẫu nhiên một học sinh trong một trường. Xét 2 biến cố sau:\\
	P: \lq\lq Học sinh đó bị cận thị\rq\rq .\\
	Q: \lq\lq Học sinh đó học giỏi Toán\rq\rq .\\
	Nêu nội dung của các biến cố P giao Q là 
	\choice
	{Biến cố PQ xảy ra khi học sinh đó chỉ bị cận thị hoặc chỉ học giỏi môn Toán}
	{Biến cố PQ xảy ra khi học sinh đó bị cận thị hoặc học giỏi môn Toán}
	{\True Biến cố PQ xảy ra khi học sinh đó vừa bị cận thị vừa học giỏi môn Toán}
	{Biến cố PQ xảy ra khi học sinh đó không bị cận thị và không học giỏi môn Toán}
	\loigiai{
		Dựa vào định nghĩa.}
\end{ex}
\begin{ex}%Câu 7.%[1D9N2-1]
	Hai bạn Dương và Khoa, mỗi người gieo đồng thời một con xúc xắc cân đối, đồng chất. Xét hai biến cố sau:\\
	E: \lq\lq Số chấm xuất hiện trên con xúc xắc bạn Dương gieo là số chính phương\rq\rq .\\
	B: \lq\lq Số chấm xuất hiện trên con xúc xắc của bạn Khoa là số chia hết cho 4\rq\rq .\\
	Chọn câu đúng: 
	\choice
	{\True Hai biên cố E và B độc lập}
	{Hai biến cố E và B không độc lập}
	{Hai biến cố E và B xung khắc}
	{Hai biến cố E và B không xung khắc}
	\loigiai{
		Dựa vào định nghĩa}
\end{ex}
\begin{ex}%Câu 8.%[1D9N1-1]
	Tung một đồng xu cân đối và đồng chất hai lần liên tiếp. Xét các biến cố:\\
	X: \lq\lq Đồng xu xuất hiện mặt sấp ở lần gieo thứ nhất\rq\rq .\\
	Y: \lq\lq Đồng xu xuất hiện mặt ngửa ở lần gieo thứ hai\rq\rq .\\
	Hai biến cố X và Y là hai 
	\choice
	{biến cố đối}
	{\True biến cố độc lập}
	{biến cố giao}
	{biến cố hợp}
	\loigiai{
		Dựa vào định nghĩa.}
\end{ex}
\begin{ex}%Câu 9.%[1D9N1-2]
	Một hộp đựng 25 tấm thẻ cùng loại được đánh số từ 1 đến 25. Rút ngẫu nhiên một tấm thẻ trong hộp. Xét các biến số P: \lq\lq Số ghi trên tấm thẻ là số chia hết cho 4\rq\rq \, và Q: \lq\lq Số ghi trên tấm thẻ là số chia hết cho 6\rq\rq.  Nội dung của biến cố giao PQ là gì?
	\choice
	{$\left\{3;6;9;12;15;18;24\right\}$}
	{$\{1;2;4\}$}
	{$\left\{12;16;20;24\right\}$}
	{\True $\{12;24\}$}
	\loigiai{
		Bội chung của $4$ và $6$ là $\left\{12;24\right\}$.}
\end{ex}
\begin{ex}%Câu 10.%[1D9N1-1]
	Một hộp có 52 chiếc thẻ cùng loại, mỗi thẻ được ghi một trong các số tự nhiên, hai thẻ khác nhau thì ghi hai số khác nhau. Rút ngẫu nhiên 1 chiếc thẻ trong hộp. Xét biến cố A: \lq\lq Số xuất hiện trên thẻ được rút ra là số lẻ\rq\rq\,và biến cố B: \lq\lq Số xuất hiện trên thẻ được rút ra là số nguyên tố\rq\rq.\\
	Chọn câu \textbf{sai} khi phát biểu về biến cố AB. 
	\choice
	{\True \lq\lq Số xuất hiện trên thẻ là số lẻ hoặc số nguyên tố\rq\rq}
	{\lq\lq Số xuất hiện trên thẻ là số nguyên tố lẻ\rq\rq }
	{\lq\lq Số xuất hiện trên thẻ là số nguyên tố khác $2$\rq\rq }
	{\lq\lq Số xuất hiện trên thẻ vừa là số lẻ vừa là số nguyên tố\rq\rq}
	\loigiai{
		Dựa vào định nghĩa có \lq\lq Số xuất hiện trên thẻ là số lẻ hoặc số nguyên tố\rq\rq\, là phát biểu sai về AB.}
\end{ex}
\begin{ex}%Câu 11.%[1D9N1-1]
	Một hộp đựng 5 viên bi màu đỏ và 6 viên bi màu xanh có cùng kích thước và khối lượng. Bạn Hưng lấy ngẫu nhiên một viên bi và không trả lại vào hộp. Tiếp theo bạn An lấy ngẫu nhiên một viên bi từ hộp đó. Xét biến cố A: \lq\lq Hưng lấy được viên bi màu đỏ\rq\rq , biến cố B: \lq\lq An lấy được viên bi màu xanh\rq\rq . Hai biến cố A và B là hai  
	\choice
	{biến cố đối}
	{biến cố độc lập}
	{\True biến cố không độc lập}
	{biến cố hợp}
	\loigiai{
		Hai biến cố A và B là hai  biến cố không độc lập vì xác suất lấy dược viên bi xanh của An phụ thuộc vào xác suất bạn Hưng lấy được bi màu nào.}
\end{ex}
\begin{ex}%Câu 12.%[1D9N2-2]
	Một tổ trong lớp 12CT có 4 học sinh nữ là Giang, Thanh, Thảo, Anh và 5 học sinh nam là Sơn, Tùng, Hoàng, Tiến, Hải. Trong giờ học, giáo viên chọn ngẫu nhiên một học sinh trong tổ đó lên bảng để kiểm tra bài. Xét các biến cố sau:\\
	H: \lq\lq Học sinh đó là một bạn nữ\rq\rq .\\
	K: \lq\lq Học sinh đó có tên bắt đầu bằng chữ T\rq\rq .\\
	Nêu nội dung của biến cố hợp H hợp K.
	\choice
	{\lq\lq Giang\rq\rq}
	{\lq\lq Tùng, Tiến, Thảo\rq\rq}
	{\True \lq\lq Tùng, Tiến, Giang, Anh, Thanh, Thảo\rq\rq}
	{\lq\lq Giang, Thanh, Thảo, Anh\rq\rq}
	\loigiai{
		Ta có H hợp K gồm các phần tử \lq\lq Tùng, Tiến, Giang, Anh, Thanh, Thảo\rq\rq.}
\end{ex}
\Closesolutionfile{ans}
\indapan{6}{ans/ans-OTC8-TN}

\cauds
\Opensolutionfile{ans}[ans/ans-OTC8-DS]
\setcounter{ex}{0}
\begin{ex}%Câu 1.%[1D9H1-3]
	Cho $A$ và $B$ là hai biến cố độc lập với nhau, biết $\mathrm{P}(A)=0{,}2$; $\mathrm{P}(B)=0{,}3$. Khi đó
	\choiceTF	
	{\True $\mathrm{P}(AB)=0{,}06$}
	{$\mathrm{P}(A\overline{B})=0{,}12$}
	{\True $\mathrm{P}(\overline{A}\overline{B})=0{,}56$}
	{\True $\mathrm{P}(\overline{A}B)=0{,}24$}
\loigiai{
	Vì $A$, $B$ là hai biến cố độc lập nên $A$, $B$, $\overline{A}$, $\overline{B}$ là các biến cố đôi một độc lập.\\
	Ta có
	\begin{itemchoice}
	\itemch Đúng. $\mathrm{P}(AB)=\mathrm{P}(A)\cdot\mathrm{P}(B)=0{,}2\cdot 0{,}3=0{,}06$.
	\itemch Sai. $\mathrm{P}(A\overline{B})=\mathrm{P}(A)\cdot\mathrm{P}(\overline{B})=0{,}2\cdot 0{,}7=0{,}14$.
	\itemch Đúng. $\mathrm{P}(\overline{A}\overline{B}) =\mathrm{P}(\overline{A})\cdot\mathrm{P}(\overline{B})=0{,}8\cdot 0{,}7=0{,}56$.
	\itemch Đúng. $\mathrm{P}(\overline{A}B)
	=\mathrm{P}(\overline{A})\cdot\mathrm{P}(B)=0{,}8\cdot 0{,}3=0{,}24$.
	\end{itemchoice}
	}
\end{ex}
	\begin{ex}%Câu 2.%[1D9H1-3]
Cho $A$, $B$ là hai biến cố độc lập và $\mathrm{P}(A)=\dfrac{1}{4}$, $\mathrm{P}(B)=\dfrac{1}{3}$. Khi đó
\choiceTF
{$\mathrm{P}(AB)=\dfrac{1}{2}$}
{$\mathrm{P}(A\overline{B})=\dfrac{1}{16}$}
{\True $\mathrm{P}(\overline{A}\overline{B})=\dfrac{1}{2}$}
{\True $\mathrm{P}(\overline{A}B)=\dfrac{1}{4}$}
\loigiai{
Ta có xác suất các biến cố đối: $\mathrm{P}(\overline{A})=\dfrac{3}{4}$, $\mathrm{P}(\overline{B})=\dfrac{2}{3}$.\\
Vì $A,B$ là hai biến cố độc lập nên mỗi cặp biến cố lấy ra từ $A$, $B$, $\overline{A}$, $\overline{B}$ đều là cặp biến cố độc lập. Do vậy
\begin{itemchoice}
	\itemch Sai. 
	$\mathrm{P}(AB)=\mathrm{P}(A)\cdot\mathrm{P}(B) =\dfrac{1}{4}\cdot\dfrac{1}{3}=\dfrac{1}{12}$.
	\itemch Sai.  $\mathrm{P}(\overline{A}B)=\mathrm{P}(\overline{A})\cdot\mathrm{P}(B)
	=\dfrac{3}{4}\cdot\dfrac{1}{3}=\dfrac{1}{4}$.
	\itemch Đúng.
	$\mathrm{P}(A\overline{B})=\mathrm{P}(A)\cdot\mathrm{P}(\overline{B})
	=\dfrac{1}{4}\cdot\dfrac{2}{3}=\dfrac{1}{6}$.
	\itemch Đúng. $\mathrm{P}(\overline{A}\overline{B})
	=\mathrm{P}(\overline{A})\cdot\mathrm{P}(\overline{B})=\dfrac{3}{4}\cdot\dfrac{2}{3}=\dfrac{1}{2}$.	
\end{itemchoice}}
\end{ex}
\begin{ex}%Câu 3.%[1D9V1-3]
Một người vừa gieo một con xúc xắc để ghi lại số chấm xuất hiện, sau đó người này tiếp tục chọn ngẫu nhiên một lá bài từ bộ bài $52$ lá. 
\choiceTF	
{\True Gọi $A$ là biến cố: \lq\lq Số chấm của xúc xắc lớn nhất\rq\rq, khi đó: $\mathrm{P}(A)=\dfrac{1}{6}$}
{\True Gọi $B$ là biến cố: \lq\lq Chọn được một lá bài tây\rq\rq, khi đó: $\mathrm{P}(B)=\dfrac{3}{13}$}
{\True Xác suất để số chấm trên con xúc xắc là lớn nhất và chọn được một lá bài tây bằng $\dfrac{1}{26}$}
{Xác suất để số chấm trên con xúc xắc và số của lá bài là giống nhau bằng $\dfrac{1}{16}$}
\loigiai{
Gọi $A$ là biến cố: \lq\lq Số chấm của xúc xắc lớn nhất\rq\rq, $B$ là biến cố: \lq\lq Chọn được một lá bài tây\rq\rq. Dễ thấy $A$, $B$ là hai biến cố độc lập.
\begin{itemchoice}
\itemch Đúng. Ta có: $A=\{6\}\Rightarrow n(A)=1\Rightarrow\mathrm{P}(A)=\dfrac{1}{6}$.
\itemch Đúng. Ta biết bộ bài 52 lá thì có 12 lá bài tây, nên xác suất chọn được một lá bài tây là $\mathrm{P}(B)=\dfrac{3}{13}$.
\itemch Đúng. Suy ra $\mathrm{P}(AB)=\mathrm{P}(A)\cdot\mathrm{P}(B)=\dfrac{1}{6}\cdot\dfrac{3}{13}=\dfrac{1}{26}$.
\itemch Sai.
Để thu được số chấm trên con xúc xắc và số của lá bài giống nhau thì ta có 6 cách để có được số chấm một con xúc xắc, ứng với mỗi cách đó thì có đúng 4 cách tìm được lá bài thoả mãn.
Việc gieo xúc xắc và rút ngẫu nhiên lá bài là độc lập.\\
Gọi $X$ là biến cố cần tính xác suất, ta có: $\mathrm{P}(X)=\dfrac{6}{6}\cdot\dfrac{4}{52}=\dfrac{1}{13}$.
\end{itemchoice}}
\end{ex}
\begin{ex}%Câu 4.%[1D9V2-3]
Trên một bảng quảng cáo, người ta mắc hai hệ thống bóng đèn. Hệ thống I gồm 2 bóng mắc nối tiếp, hệ thống II gồm 2 bóng mắc song song. Khả năng bị hỏng của mỗi bóng đèn sau 6 giờ thắp sáng liên tục là 0{,}15. Biết tình trạng của mỗi bóng đèn là độc lập. Khi đó xác suất để
\choiceTF
{\True Hệ thống II bị hỏng (không sáng) bằng $0{,}0225$}
{\True Hệ thống II hoạt động bình thường bằng $0{,}9775$}
{Hệ thống I bị hỏng (không sáng) bằng $0{,}5775$}
{Cả hai hệ thống bị hỏng (không sáng) bằng $\approx 0{,}02624$}
\loigiai{
	\begin{itemchoice}
\itemch Đúng. Nhận xét: Hệ thống II gồm 2 bóng được mắc song song nên nó chỉ hỏng khi cả hai bóng đều hỏng.\\
Gọi $B$ là biến cố: \lq\lq Hệ thống II bị hỏng\rq\rq, ta có: $\mathrm{P}(B)=0{,}15\cdot 0{,}15=0{,}0225$.
\itemch Đúng. Xác suất để hệ thống II hoạt động bình thường: $\mathrm{P}(\overline{B})=1-0{,}0225=0{,}9775$.
\itemch Sai. Nhận xét: Hệ thống I chỉ hoạt động bình thường khi cả hai bóng bình thường.\\
Gọi $A$ là biến cố: "Hệ thống I bị hỏng". Khi đó xác suất để hệ thống I hoạt động bình thường là $\mathrm{P}(\overline{A})=0{,}85\cdot 0{,}85=0{,}7225$.\\
Suy ra $\mathrm{P}(A)=1-0{,}7225=0{,}2775$.
\itemch Sai. Xác suất để cả hai hệ thống I, II đều bị hỏng là\\
$\mathrm{P}(AB)=\mathrm{P}(A)\cdot\mathrm{P}(B)=0{,}2775\cdot 0{,}0225=\dfrac{999}{160000}\approx 0{,}00624$.
\end{itemchoice}
}
\end{ex}
\Closesolutionfile{ans}
\indapan{3}{ans/ans-OTC8-DS}

\Opensolutionfile{ans}[ans/ans-OTC8-KQ]
\caukq
\setcounter{ex}{0}
\begin{ex}%Câu 1.%[1D9V1-3]
	Một bình đựng $5$ viên bi xanh và $3$ viên bi đỏ (các viên bi chỉ khác nhau về màu sắc). Lấy ngẫu nhiên một viên bi, rồi lấy ngẫu nhiên một viên bi nữa. Tính xác suất của biến cố \lq\lq Lấy lần thứ hai được một viên bi xanh\rq\rq\, (làm tròn kết quả đến hàng phần trăm).
	\shortans{$0{,}63$}
\loigiai{
	Gọi $A$ là biến cố \lq\lq Lấy lần thứ hai được một viên bi xanh\rq\rq .\\
	Gọi $B$ là biến cố \lq\lq Lấy lần thứ nhất được bi xanh, lấy lần thứ hai cũng được một bi xanh\rq\rq .\\
	Xác suất biến cố $B$ là $\mathrm{P}(B)=\dfrac{5}{8}\cdot\dfrac{4}{7}=\dfrac{5}{14}$.\\
	Gọi $C$ là biến cố \lq\lq Lấy lần thứ nhất được bi đỏ, lấy lần thứ hai được bi xanh\rq\rq .\\
	Xác suất biến cố $C$ là $\mathrm{P}(C)=\dfrac{3}{8}\cdot\dfrac{5}{7}=\dfrac{15}{56}$.\\
	Ta có $A=B\cup C$, vì $B$, $C$ là hai biến cố xung khắc nên ta có\\
	$\mathrm{P}(A)=\mathrm{P}(B\cup C)=\mathrm{P}(B)+\mathrm{P}(C)=\dfrac{5}{14}+\dfrac{15}{56}=\dfrac{35}{56}=\dfrac{5}{8}$.}
	\end{ex}
	\begin{ex}%Câu 2.
Một chiếc máy có 2 động cơ I và II hoạt động độc lập với nhau. Xác suất để động cơ I chạy tốt và động cơ II chạy tốt lần lượt là $0{,}9$ và $0{,}8$. Tính xác suất để có ít nhất 1 động cơ chạy tốt (làm tròn kết quả đến hàng phần trăm).
\shortans{$0{,}98$}
\loigiai{
Gọi $A_i$ là động cơ thứ $i$ chạy tốt.\\
Gọi $A$ là biến cố \lq\lq  có ít nhất một động cơ chạy tốt\rq\rq .\\
$\overline{A}$ là biến cố \lq\lq  không động cơ nào chạy tốt\rq\rq .\\
Ta có: $\overline{A}=\overline{A_1}\overline{A_2}\Rightarrow\mathrm{P}\left(\overline{A}\right)=\mathrm{P}\left(\overline{A_1}\right)\mathrm{P}\left(\overline{A_2}\right)=(1-0{,}9)(1-0{,}8)=0{,}02$.\\
Vậy $\mathrm{P}(A)=1-\mathrm{P}\left(\overline{A}\right)=0{,}98$.}
\end{ex}
\begin{ex}%Câu 3.%[1D9V2-3]
Một hộp chứa $5$ viên bi đỏ, $6$ viên bi xanh và $7$ viên bi vàng (các viên bi kích thước như nhau). Lấy ngẫu nhiên $4$ viên bi từ hộp. Tính xác suất để trong $4$ viên bi lấy được có nhiều nhất 2 viên bi đỏ (làm tròn kết quả đến hàng phần trăm).
\shortans{$0{,}39$}
\loigiai{
Lấy ngẫu nhiên $4$ viên bi từ hộp. Số kết quả có thể xảy ra là $n({\Omega}_1)=\mathrm{C}_{18}^4$.\\
Gọi $A$ là biến cố: \lq\lq  $4$ viên bi lấy được trong hộp có nhiều nhất 2 viên bi đỏ\rq\rq .\begin{itemize}
\item[TH1:] Không lấy được viên bi đỏ nào trong $4$ bi có $\mathrm{C}_5^{0}\cdot\mathrm{C}_{13}^4$ cách.
\item[TH2:] Lấy được $1$ bi đỏ trong $4$ bi có $\mathrm{C}_5^1\cdot\mathrm{C}_{13}^3$ cách.
\item[TH3:] Lấy được $2$ bi đỏ trong $4$ bi có $\mathrm{C}_5^2\cdot\mathrm{C}_{13}^2$ cách.
\end{itemize}
Số kết quả thuận lợi cho $A$ là
$n(A)=\mathrm{C}_5^0\cdot\mathrm{C}_{13}^4+\mathrm{C}_5^1\cdot\mathrm{C}_{13}^3+\mathrm{C}_5^2\cdot\mathrm{C}_{13}^2$.\\
Khi đó, xác suất để trong $4$ viên bi lấy được trong hộp có nhiều nhất 2 viên bi đỏ là\\
$\mathrm{P}(A)=\dfrac{n(A)}{n({\Omega}_1)}=\dfrac{\mathrm{C}_5^0\cdot\mathrm{C}_{13}^4+\mathrm{C}_5^1\cdot\mathrm{C}_{13}^3+\mathrm{C}_5^2\cdot\mathrm{C}_{13}^2}{\mathrm{C}_{18}^4}=\dfrac{481}{1224}=0{,}39$.}
\end{ex}
\begin{ex}%Câu 4.%[1D9V2-3]
Một chiếc hộp có chín thẻ đánh số thứ tự từ $1$ đến $9$. Rút ngẫu nhiên $2$ thẻ rồi nhân hai số ghi trên thẻ lại với nhau. Tính xác suất để kết quả nhân được là một số chẵn (kết quả làm tròn đến hàng phần trăm).
\shortans{$0{,}72$}
\loigiai{
Trường hợp 1: hai số rút ra đều là số chẵn: $P_1=\dfrac{\mathrm{C}_4^2}{\mathrm{C}_9^2}=\dfrac{1}{6}$.\\
Trường hợp 2: hai số rút ra có một số lẻ, một số chẵn: $P_2=\dfrac{\mathrm{C}_4^1\cdot\mathrm{C}_5^1}{\mathrm{C}_9^2}=\dfrac{5}{9}$.\\
Vậy xác suất để kết quả nhân được là một số chẵn là $P=P_1+P_2=\dfrac{1}{6}+\dfrac{5}{9}=\dfrac{13}{18}$.}
\end{ex}
\begin{ex}%Câu 5.%[1D9V1-3]
Gieo đồng xu đồng chất $3$ lần liên tiếp. Tính xác suất để cả $3$ lần đều được mặt sấp  (kết quả làm tròn đến hàng phần trăm).
\shortans{$0{,}13$}
\loigiai{
\textbf{Cách 1.} Gọi $A_1\colon$\lq\lq Gieo lần đầu mặt sấp\rq\rq  $\Rightarrow\overline{A_1}\colon$\lq\lq Gieo lần đầu mặt ngửa\rq\rq.\\
$A_2\colon$\lq\lq Gieo lần 2 mặt sấp\rq\rq  $\Rightarrow\overline{A_2}\colon$\lq\lq Gieo lần 2 mặt ngửa\rq\rq .\\
$A_3\colon$\lq\lq Gieo lần 3 mặt sấp\rq\rq  $\Rightarrow\overline{A_3}\colon$\lq\lq Gieo lần 3 mặt ngửa\rq\rq.\\
$ \Rightarrow\mathrm{P}(A_1)=\mathrm{P}(A_2)=\mathrm{P}(A_3)=\mathrm{P}\left(\overline{A_1}\right)=\mathrm{P}\left(\overline{A_2}\right)=\mathrm{P}\left(\overline{A_3}\right)=0{,}5 $.\\
Xác suất để cả 3 lần gieo là sấp $0{,}5\cdot 0{,}5\cdot 0{,}5=0{,}125$.\\
\textbf{Cách 2: }Dùng sơ đồ hình cây: 
% TODO: \usepackage{graphicx} required
\begin{center}
	\includegraphics[scale=0.6]{image/Otc8Cau5TLN}
\end{center}
Dựa và sơ đồ ta thấy:\\
Số phần tử không gian mẫu là $n(\Omega)=8$.\\
Số phần tử của biến cố là $n(A)=1$.\\
Xác suất biến cố là $\mathrm{P}(A)=\dfrac{1}{8}=0{,}125$.}
\end{ex}
\begin{ex}%Câu 6.%[1D9V1-3]
Gieo hai con súc sắc I và II cân đối, đồng chất một cách độc lập. Ta có biến cố $A$: \lq\lq Có ít nhất một con súc sắc xuất hiện mặt $6$ chấm\rq\rq . Lúc này giá trị của $\mathrm{P}(A)$ là
 (kết quả làm tròn đến hàng phần trăm).
\shortans{$0{,}31$}
\loigiai{
Gọi $A_i(i=1;2)$ là biến cố: \lq\lq Con súc sắc thứ $i$ ra mặt $6$ chấm\rq\rq \\
$ \Rightarrow A_1 $ và $A_2$ là hai biến cố độc lập và ta có $\heva{&\mathrm{P}(A_1)=\dfrac{1}{6}\\&\mathrm{P}(A_2)=\dfrac{1}{6}.}$ \\
Thay vì tính $\mathrm{P}(A)$ ta đi tính $\mathrm{P}\left(\overline{A}\right)$. Ta có $\overline{A}=\overline{A_1}\cdot\overline{A_2}$.\\
$\mathrm{P}\left(\overline{A}\right)=\mathrm{P}\left(\overline{A_1}\right)\cdot\mathrm{P}\left(\overline{A_2}\right)=\left(1-\mathrm{P}(A_1)\right)\cdot\left(1-\mathrm{P}(A_2)\right)=\dfrac{5}{6}\cdot\dfrac{5}{6}=\dfrac{25}{36}$.\\
Vậy $\mathrm{P}(A)=1-\mathrm{P}\left(\overline{A}\right)=1-\dfrac{25}{36}=\dfrac{11}{36}$.}
\end{ex}
\Closesolutionfile{ans}
\indapan{6}{ans/ans-OTC8-KQ}